\begin{frame}
	\frametitle{Ответы на замечания Брантова~А.В.}
	\textbf{Замечание}: Во второй главе диссертационной работы при подавлении взаимного влияния одномодового и межмодового четырехволнового смещений наблюдается возбуждение ВКР сигнала. Было бы полезно описать процесс возникновения ВКР более подробно с указанием условий фазового синхронизма (частоты оптических фононов в рассматриваемом оптоволокне) и оценки порога возникновения ВКР. 
		
	\textbf{Ответ:} Для германосиликатного ОВ спектр усиления ВКР имеет широкий максимум со сдвигом 13--14~ТГц, что соответствует частоте оптических фононов. При накачке 1030~нм это даёт стоксову компоненту вблизи 1080~нм, которая и наблюдалась экспериментально (рис.~2.16б). Условие фазового синхронизма при ВКР выполняется автоматически: частота стоксовой волны определяется частотами волны накачки и оптических фононов, поэтому специального согласования, в отличие от ЧВС, не требуется. Спектральные оценки показывает, что при длине ОВ порядка 4~м и пиковой мощности накачки 4~кВт в отсутствие ЧВС доля ВКР составляет порядка 10\%
	
\end{frame}


\begin{frame}
\frametitle{Ответы на замечания Брантова~А.В.}
	
	\textbf{Замечание}: В третьей главе приводится несимметричная картина профиля пучка на второй гармонике (Рис. 3.3). С чем может быть связана возникающая асимметрия?  
	
	\textbf{Ответ:} Наблюдаемая асимметрия профиля пучка второй гармоники имеет две составляющие. Во-первых, она частично обусловлена модовым составом излучения на длине волны 515~нм: волокно FUD-3562 на этой длине волны поддерживает до девяти поперечных мод, и условие фазового синхронизма для ГВГ может выполняться не только для перехода в моду LP$_{12}$. Во-вторых, часть наблюдаемой асимметрии связана с условиями регистрации: изображение на рис.~3.3 получено КМОП-камерой без дополнительной юстировки, поскольку её основной целью была качественная визуализация неодномодовости пучка.	
\end{frame}


\begin{frame}
	\frametitle{Ответы на замечания Брантова~А.В.}
	\textbf{Замечание}: Во второй главе проводится моделирование генерации излучения в оптоволокне на второй гармонике с использованием модельных уравнений с постоянными константами и отмечается существенное отличие теоретических результатов от экспериментальных на этапе формирования решетки, что вероятно связано с временной зависимостью используемых в модели констант.  Было бы полезно продемонстрировать, что введение временной зависимости позволяет приблизить теоретические результаты к экспериментальным и изменить скорость эффективности генерации второй гармоники (Рис. 3.4)

	
	\textbf{Замечание}: Чтение диссертации затрудняет введение большого числа сокращений, не все из которых приведены в списке сокращений  

\end{frame}